\section{Reactionary Chic}

There was a time --- not so distant in terms of world history --- when converts came from world-class authors, novelists, or philosophers. I am thinking in particular of Jaris-Karl Huysmans, Charles Peguy, Leon Bloy, Jacques Maritain, G K Chesterton, or Graham Greene, not to mention Charles Maurras and Ernst Junger, all exemplars of the Catholic intellectual.

In our time, that is not the case. It seems that in the Internet age, there is an abundance of converts (or just as many so-called ``reverts'') who feel compelled to write and speak publicly via blogs, youtube, and even satellite radio. Although I am sure that the angels in heaven are rejoicing over their conversions, on the human plane we might hope that they take a vow of silence instead.

When they are not being progressive --- demonstrating that their conversion is still incomplete --- they too often become too negative, not about the worldview they had rejected, but rather about their newly adopted Tradition. It is as though the solidity and certainty they were expecting were ephemeral.

The memorization of the catechism is an insufficient compared to a solid study of Tradition. These new converts can bowl you over with knowledge of dogma, but what is really needful is a solid grounding in Tradition. Hence, there is little added in terms of art, literature, and so on. That is the point of \textbf{Ananda Coomaraswamy}'s observation:

\begin{quotex}
A knowledge of modern Christianity will be of little use because the fundamental sentimentality of our times has diminished what was once an intellectual doctrine to a mere morality that can hardly be distinguished from a pragmatic humanism.
\end{quotex}


When the intellectual tradition was better known, artists and philosophers could often count on support from wealthy patrons. But as for ``mere morality'', there are better alternatives today. A lukewarm leftism is the null hypothesis in our time, which therefore is difficult to displace. A convert has little to add to contemporary morality, other than perhaps some quaint objections to abortion, divorce, or homosexuality, but the force of those counter-cultural positions is diminishing.

\paragraph{Making Reaction Chic Again}


Wealthy young patrons like the chic Larry and Toby\footnote{\url{https://www.nytimes.com/2017/06/24/fashion/larry-toby-milstein-leonard-bernstein-ghost.html}} support secular hospitals and art. So how to make reaction chic? Instead of séances to mediocre historical figures, how about invocations to the saints and the Holy Guardian Angel? Then there can be meditation classes, like those described in these pages. Kabbalah is somewhat chic. But we could offer Hermetic and Tarot centres. Fundamentally, the Western Tradition is so alien to the modern mind, it could become chic again just out of sheer novelty. The goal would be to reach the TOOS and Category X status categories. That would require an elite view, something more elite than the genteel leftism that passes for elitism. So that leads to the next topic.

\paragraph{The Futility of Intelligence}


This study, The Inappropriately Excluded\footnote{\url{http://polymatharchives.blogspot.co.uk/2015/01/the-inappropriately-excluded.html}}, shows that those with an ultra-high IQ (UHIQ) are nearly systematically excluded from many of the professions which are normally associated with high intelligence. The sweet spot of IQ seems to be in the 120-135 range (very high IQ -- VHIQ). Other than theoretical physics and advanced mathematics, all college majors can be mastered by those in that range. Mastery for the VHIQ is the ability to remember and repeat the material that has been taught. All their lives, they have been assured about their intelligence, and therefore see little need to doubt it. They don't suffer from the adjustment problems that often plague the UHIQ.

The UHIQ not only masters the material, but goes beyond it. In some cases, they may then pester the teachers about their mistakes. In other cases, they find the behaviour of the less intelligent so inexplicable, that they have trouble adjusting to classmates. Without getting bogged down in this topic, the point is that the UHIQ for the most part absorb leftism in the belief that is what intelligent people accept. An alternative viewpoint might be coherently proposed that could perhaps gain some traction. I'm afraid that conventional conservatism is too prole and even incoherent. The alt-right is too vulgar and lacks moral imagination. Perhaps neo-reaction aspires to higher intellectual standards, but it has not gained traction.

\paragraph{Symbolism and Realism}


In our age, the ideal of ``realism'' in art is considered edgy. The problem is that the notion of what is ``real'' is so constricted. For example, I recently caught a fragment of a TV show (I forget which one), in which a chic lesbian high school student was discussing her class text in which ``masturbation'' was mentioned in a positive way. I assume it was not one of the Great Books, even though Thomas Aquinas did address the question from a different perspective.

The opposite is Symbolism, which uses symbols and imagination to reach deeper dimensions of reality. There are soft echoes of that today, although the themes are attenuated. A recent example is the return of \emph{Twin Peaks}, which is inspiring a lot of speculation about its ``deeper'' meaning. Unfortunately, there simply is no deeper meaning, as though the dream sequences were sufficient in themselves.

The point is that Art is the path to Chic, and the new converts are without Art.

\paragraph{Free Will and Animals}


In \emph{The Romance of the Rose}, there is speculation about what would happen if animals had free will. Perhaps the hens would refuse to lay eggs, the cattle would conspire against the cowboys, and your dog would stop sniffing your ass. Of course, that is exactly what human life is like, thus proving free will.

\paragraph{Forced Conversions}


Nevertheless, freedom is nauseating, so it is often denied. Men are happiest when they can attribute their behaviour to external causes, e.g., genetics, a disease, environment or circumstances.

A logical argument is likewise a form of compulsion. Recently, I've heard two different ``explanations'' from otherwise intelligent men about why they apostatised from the religion of their ancestors. One mentioned that he recognized the pagan roots of Christianity and the other, after he read some historian's claim that Christianity was reducible to the ancient Egyptian religion.

Quite strange, since these pages have often documented those same points. Hence, the arguments are without force, and indicate bad faith since they deny free will. A conversion is a matter of will, you will it, or you will something else.

\paragraph{Birth and Irrationality}


The Hermetic teaching, as revealed by Valentin Tomberg is this:

\begin{quotex}
The individuality descends consciously and of his own free will to birth, into an environment where he is wanted and awaited.
\end{quotex}


What is remarkable is that some Russian theologians like \textbf{Vladimir Solovyov} and \textbf{Sergius Bulgakov} integrated some of these teachings into their systems. In \emph{The Burning Bush}, Bulgakov discusses the problem of births. He begins with this:

\begin{quotex}
God sends souls into bodies individually fitted with qualities ... all the circumstances of birth and earthly destiny.
\end{quotex}


However, the soul cannot be free if it is created against its own will. Hence, there is an element of self-determination in the soul when it is created. The soul is joined to its body \emph{willingly}. The obvious question is why does not the soul remember having made this choice. The answer is that it is not a temporal event in the life of the soul, but is pre-temporal.

The point is that the human spirit assumes a personality with individual qualities; men are born different and have different destinies. These qualities are partially determined by the individual heredity which includes family, clan, and even its generation.

If you observe yourself carefully, you may notice little clues that help you recall your primordial choice. Perhaps some bon mot that means little to the world, yet something special to your own life.

\paragraph{Skin in the Game}


\textbf{Nassim Taleb} explains the necessity of eating one's own turtles\footnote{\url{https://medium.com/incerto/why-each-one-should-eat-his-own-turtles-revised-8a4be2f11e61}}. More prosaically, this may be termed ``practice what you preach''. Yes, it is true that I drink wine, cook pasta, and read medieval philosophy while listening to Gregorian chant, or to Erik Satie when I am feeling ``edgy''.

Mr. Taleb has convinced me that the personal element is relevant to the intellectual element. This is unlike previous generations of writers on Tradition who avoided any hint of the personal. However, in the course of life, I've encountered False Seeming much too often; the initiate must learn to distinguish what is from what seems to be.

I don't believe that so-called white nationalists have been forthcoming with their ancestry reports. I'm considering posting mine publicly, if I determine that it is a wise or worthwhile idea. Basically, it is 97\% European, with 2.5\% Middle Eastern. Of course, that is, or was, part of the Roman Empire which is my real identity, although there is no big market for Romanity at the moment. Nevertheless, it probably accounts for my fondness for falafel and belly dancers.

\paragraph{Secret Rebel}


You don't need extravagant outward displays to be convincing as a societal rebel. In the spirit of openness, here is a fragment from my personality profile:

\begin{quotex}
He has a good deal less than optimal concern for the standards of society and tends to be low on conventionality and group conformity.
\end{quotex}


The psychologist offered to cure me of that, but I graciously declined.

\paragraph{Calumny and Detraction}


The initiate should be very careful not to repeat things about other people which should not be told. Nevertheless, Evil Tongue is unable to be quiet. Now calumny, which is the making of false and defamatory statements, is commonly acknowledged to be a moral flaw. At least it used to be before political positions became warlike.

What is not as well-known is that detraction is also immoral. In other words, truth is no defence. Detraction is the making of true and defamatory statements about someone's flaws, if they are not publicly known and the only reason for their revelation is to hurt someone. Evil Tongue simply cannot help himself.

\paragraph{A Child's Prayer}


When I was a young boy, I prayed this prayer daily.

\begin{itemize}


\item Thank you, God, for making me Catholic.

\item Thank you, God, for making me an American.

\item Thank you, God, for making me Italian
\end{itemize}


Point 1 is obvious, since everyone else is risking his salvation.

The reason for Point 2 may not be. At that time, the USA was locked in a cold war with the Russians, who were promoting communism and atheism. Those went against point 1. At this point in history, there is some ambiguity since the roles of the two parties seem to be reversing.

Point 3 is because we had more fun than the other kids and ate much better. Once time I was invited to a friend's house for dinner. The mother must have thought she was pleasing me when she served ravioli from a can. I ran out of the house and back home when no one was looking. They must have had a panic attack, but they eventually called my mom.


\flrightit{Posted on 2017-07-10 by Cologero}

\begin{center}* * *\end{center}

\begin{footnotesize}
\begin{sffamily}
\texttt{James on 2017-07-11 at 04:17 said:}

Thanks for another great post.

\hfill

\texttt{Fred on 2017-07-11 at 06:40 said:}

The problem is can we offer and make chic a western tradition we lost?

There is hardly any 19th or 20th century hermeticist who didn't use karma and other eastern concepts to explain his theories. The presence of reincarnation is probably the easiest way to spot the frauds.

The only time I felt part of a living tradition is when I studied for confirmation with Catholic traditionalists: thomism, Aristotle, the Pseudo-Dyonisius, Plato. Solid foundations. The only time.

Buddhism is chic today in the west, too bad in the west is mostly leftist new age self-centered garbage.

Hinduism is chic today in the west, too bad is just misconceptions and ``do as you please''.

Why Hermeticism would be different than another parody?

\hfill

\texttt{Thomas Blanchard on 2017-07-12 at 12:14 said:}

Cologero:

Speaking of Twin Peaks, I was researching a bit about David Lynch the other day and discovered that he's a long-time advocate and devotee of Maharishi Mahesh Yogi's ``Transcendental Meditation'' technique. You've mentioned familiarity with this work multiple times in the past, but not discussed it in detail. I'm curious to know your thoughts on his work.

Lynch can be found discussing the technique and some of his related insights here: \url{https://www.youtube.com/watch?v=BH4qD5Fzyjk}. He seems sound insofar as he posits God as uncreated consciousness, also reciting some perennialist boilerplate about all traditions having access to the same metaphysical truths, etc. The most glaring flaw that I can see is his attitude toward suffering, which seems incomplete, at best (he seems to view the whole phenomenon of human suffering as totally unnecessary and easily overcome by meditation).

Going by their own descriptions (\url{https://www.davidlynchfoundation.org/about-tm.html}), TM is different from ``focused meditation'' (lectio divina, jesus prayer, etc.) and ``open monitoring'' (the kind of awareness exercises we practice), and is described vaguely as ``automatic self-transcending''. The marketing on the various TM-affiliated websites is very heavy handed and big on promising huge results for little effort. Is this kind of technique totally a scam and/or counter-initiatory, or is there something of value here worth studying?

\hfill

\texttt{Numen on 2017-07-12 at 14:01 said:}

I thought similarly when I was a wee lad. Born in the States, I went on vacation to the fatherland in Bangladesh with my family, and those were some of my earliest memories. I remember my aunt telling me not to play or talk to the ``asura boys'' since they were filthy and lower than us. This aunt is the sweetest person I know, so it may sound strange to the uninitiated that she would say such a thing. Part of Tradition is nostalgic, in the sense that Schuon and Tomberg framed it as a tend toward primordiality and the sacred. Proper focus on personal aspects, I also think, can help.

\hfill

\texttt{Fred on 2017-07-13 at 05:35 said:}

@Thomas Blanchard

I never had any actual personal experience with TM but there is a lot to be suspicious around it.

First of all the fact that its founder got very rich, always the easiest way to spot false teachers.

Then it spread mostly in the anti-traditional environment of 60s and 70s hippies, californian celebrities, and other ``do what you wish'', hedonistic, self-centered kind of people.

There is actually an article of Rama Coomaraswamy against TM, Mahesh Yogi is called a rakshasa.

It doens't mean, however, there is nothing true in the technique. He probably stole something from various yoga traditions in India.

\hfill

\texttt{Arthur Konrad on 2017-07-13 at 07:15 said:}

It had a rough start few years ago, plagued then by egoistic outbursts, but NRx seems to have solidified recently into a discourse which at the very least won't compromise eloquence and good manners. It's a sorely lacking thing

The question isn't whether promoting reactionary views is possible within modern circumstances, especially when it comes to inner life, which is crucial in reaching the proper comprehension, but whether history needs another ``uncompromising'' and ``stubborn'' Traditionalist to repeat the same motions once again -- initial enthusiasm, followed by resignation, and concluded by pessimism and abstention.

Perhaps we should be content in acknowledging that this latter phase is inherent to this day and age, and go with this preposition. I'm wondering myself.

\hfill

\texttt{Schoolboy on 2017-07-13 at 17:11 said:}

Off topic but would like some input from those higher than me. Recently I read ``Lost Horizon'' by James Hilton and will be viewing the film version on Turner Classic Movies tonight. On some level Shangri La appears to be a heaven world, being reached thru arduous effort or grace from those who can offer a great opportunity if we can make use of it. All are offered the option to stay and improve themselves in ways not always available to most people. Message: be prepared to step into that new life and move into it gracefully, after all that may be the whole point of all this endless meditation, seeking, saving the whales, etc.

\hfill

\texttt{yourusernameistakeng on 2017-07-14 at 19:14 said:}

This is great. I am a musician planning to tour and I have often throught... how could I represent Tradition in my band branding etc... ..I think there is a young generation interested in this idea, but they're definitely in the minority. If they're like me they're looking and hoping for more connections to the past in culture.

\hfill

\texttt{Fred on 2017-07-16 at 16:44 said:}

A chinese saying goes like this ``It's useless to stop an avalanche with your hands, better take shelter and wait for it to end''

Maybe this is what it needs to be done. Creating safe places of faith, tradition and knowledge like the Benedictine monks during the dark centuries between the fall of Rome and Charlemagne.

\hfill

\texttt{Max on 2017-07-25 at 09:20 said:}

From where do intelligent people get the misstaken assumption that intelligent people ascribe to leftism? Because intelligent people knows not to upset the masses in vain, they proceed to tell them what they want to hear, and missapply the imperial strategy of worshiping at the shrine of the conquered as a display of good will. Most things can be assimilated except the outright denial of truth.

Note in which way Christianity fits this ancient pattern since Christ never was nor will be conquered. Christ magnanimously adopted the worthy elements of paganism as a sign of good will towards the conquered, also being an act of free will since he was not compelled to do it. When neo-pagans missinterpret an act of free will as a display of weakness it reveals that they personally know of nothing but compulsion.

A sign of reasonably high intelligence is to begin to understand the workings of ones mind. A sign of ultra high intelligence is to also understand other minds, making intelligence universal in scope rather than individually restricted.

\end{sffamily}
\end{footnotesize}

